% ==================== 附录 ====================

\section*{附录}
\addcontentsline{toc}{section}{附录}

\subsection*{附录A\quad 支撑材料文件清单}

完整可运行程序、输入数据、输出结果和正文图片均放入电子支撑材料。表~\ref{tab:appendix-files}~给出目录级清单，具体文件名、官方链接、用途和核验状态分别记录在各问题的来源清单与随包说明中。

\begin{longtable}{p{3.0cm}p{9.3cm}}
\caption{支撑材料文件清单}\label{tab:appendix-files}\\
\toprule
目录 & 内容 \\
\midrule
\endfirsthead
\toprule
目录 & 内容 \\
\midrule
\endhead
\path{matlab/q1/} & 6个MATLAB程序、数据来源清单，以及模型参数、预测、滚动验证、敏感性和自检结果 \\
\path{matlab/q2/} & 13个MATLAB程序、输入证据，以及权重、排名、敏感性和自检结果 \\
\path{matlab/q3/} & 12个MATLAB程序，以及画像、方案、排名、敏感性和自检结果 \\
\path{matlab/q4/} & 10个MATLAB程序、1个人口数据提取Python程序，以及基准、情景、模拟、敏感性和自检结果 \\
问题二数据来源目录 & 9项学术与官方来源的本地归档和书目元数据，详见\path{q2_source_manifest.csv} \\
问题三数据来源目录 & 指南参数表、评分规则、提取记录、证据说明和来源清单；大型官方原文以链接和哈希值追溯 \\
问题四数据来源目录 & 27年人口输入、情景参数、证据说明和来源清单；大型WPP原始文件以官方链接和SHA-256追溯 \\
\path{figures/q1--q4/} & 正文采用的全部PNG图片，以及用于结果复核的补充排序图和画像图 \\
\bottomrule
\end{longtable}

\subsection*{附录B\quad 完整程序索引}

问题一的运行入口为\path{matlab/q1/main_question1.m}；问题二、问题三和问题四的运行入口依次为\path{matlab/q2/run_q2_all.m}、\path{matlab/q3/run_q3_all.m}和\path{matlab/q4/run_q4_all.m}。各入口程序调用同目录下的拟合、排序、敏感性分析、绘图和自检函数，调用顺序见入口程序文件头和随包提供的《支撑材料说明》。问题四另保留\path{matlab/q4/prepare_q4_inputs.py}，用于从WPP官方单岁人口文件生成模型输入。

问题四WPP原始压缩CSV与大型Excel合计超过20\,MB，不直接装入最终支撑包；支撑包保留由其确定性提取得到的27年最小人口CSV、提取程序、官方URL和SHA-256。论文不使用代码截图，完整源程序均以可编辑文本文件随电子支撑材料提交。
